\section{Conclusion}

GEMM-based convolution remains a popular approach due to the widespread availability and heavy optimization of BLAS libraries, despite the computational and memory overhead of the commonly used Im2col transformation.
This work introduces the \algname algorithm, which removes the need for Im2col by enabling GEMM calls to access convolution inputs directly.
\algname improves upon prior GEMM-based methods by requiring no data copies, relying on unmodified GEMM kernels, and supporting a wide range of convolution configurations.
As demonstrated in the evaluation, \algname achieves significant speedups over prior work and state-of-the-art PyTorch implementations.
When integrated into PyTorch, \algname achieves up to a 1.47$\times$ speedup in end-to-end model inference, being particularly effective in image segmentation models that utilize dilated convolution.
This work highlights the potential to leverage existing high-performance GEMM kernels in an unconventional manner: by exploiting non-contiguous matrix inputs via its \texttt{ldx} parameters.
\algname's performance is limited for depthwise and heavily grouped convolutions, where the number of GEMM calls grows with group count, reducing kernel efficiency.
The algorithm is also CPU-specific: adapting it to GPUs would require addressing the kernel-launch overhead of the many small GEMM calls that \algname issues.
The selection heuristic is simple and can be further tuned per hardware platform and BLAS library.
Future work includes adapting the algorithm to target GPU execution.